% ============================================================
%  CleverMathematics — Nuanced Analysis LaTeX Template
%  Overleaf-compatible. Compile twice with pdflatex.
%  Based on the Polynomial Analysis & Quadratics PDF examples.
% ============================================================
\documentclass[11pt,a4paper]{article}

% ── Packages ──────────────────────────────────────────────────────────────────
\usepackage[T1]{fontenc}
\usepackage[utf8]{inputenc}
\usepackage[a4paper,top=22mm,bottom=22mm,left=18mm,right=18mm,headheight=14pt]{geometry}
\usepackage{amsmath,amssymb,amsthm,mathtools}
\usepackage[dvipsnames,svgnames,table]{xcolor}
\usepackage{graphicx}
\usepackage{tikz}
\usetikzlibrary{arrows.meta}
\usepackage{array,booktabs,multirow,colortbl}
\usepackage{enumitem}
\setlist[enumerate,1]{label=\textbf{\arabic*.},leftmargin=*,itemsep=4pt}
\setlist[enumerate,2]{label=(\alph*),leftmargin=2em,itemsep=2pt}
\setlist[itemize]{leftmargin=*,itemsep=2pt}
\usepackage{framed}        % for \shaded boxes — fast, no TikZ
\usepackage{parskip}
\usepackage{fancyhdr}
\usepackage{hyperref}
\hypersetup{colorlinks=true,linkcolor=NavyBlue,urlcolor=NavyBlue}

% ── Colour Palette ────────────────────────────────────────────────────────────
\definecolor{CMBlue}{RGB}{13,71,161}
\definecolor{CMAmber}{RGB}{200,110,0}
\definecolor{CMRed}{RGB}{183,28,28}
\definecolor{CMGray}{RGB}{245,245,245}
\definecolor{CMDark}{RGB}{55,65,81}
\definecolor{BoxBlue}{RGB}{227,242,253}
\definecolor{BoxBlueBd}{RGB}{30,136,229}
\definecolor{BoxAmber}{RGB}{255,248,225}
\definecolor{BoxAmberBd}{RGB}{200,130,0}
\definecolor{BoxRed}{RGB}{255,235,238}
\definecolor{BoxRedBd}{RGB}{183,28,28}
\definecolor{BoxGreen}{RGB}{232,245,233}
\definecolor{BoxGreenBd}{RGB}{46,125,50}
\definecolor{BoxPurple}{RGB}{243,229,245}
\definecolor{BoxPurpleBd}{RGB}{106,27,154}
\definecolor{BoxYellow}{RGB}{255,253,231}

% ── Box title helper ──────────────────────────────────────────────────────────
\newcommand{\boxtitle}[2]{%
  \noindent\colorbox{#1}{\parbox{\linewidth}{\vspace{1.5pt}\hspace{5pt}%
    {\small\bfseries\color{white}#2}\hspace{5pt}\vspace{1.5pt}}}\par\nobreak}

% ── Named Box Environments (framed shaded) ────────────────────────────────────
% Blue: command term, compulsory core, micro-box
\newenvironment{ctbox}[1][Command Term Spotlight]{%
  \vspace{4pt}\boxtitle{BoxBlueBd}{#1}%
  \colorlet{shadecolor}{BoxBlue}\begin{shaded}\vspace{-2pt}}%
  {\vspace{-2pt}\end{shaded}\vspace{2pt}}

\newenvironment{microbox}[1][What you need to start this Part]{%
  \vspace{4pt}\boxtitle{BoxBlueBd!80}{#1}%
  \colorlet{shadecolor}{BoxBlue!40}\begin{shaded}\vspace{-2pt}}%
  {\vspace{-2pt}\end{shaded}\vspace{2pt}}

% Amber: working required / first line given
\newenvironment{amberbox}[1][]{%
  \vspace{4pt}%
  \ifx&#1&\else\boxtitle{BoxAmberBd}{#1}\fi%
  \colorlet{shadecolor}{BoxAmber}\begin{shaded}\vspace{-2pt}}%
  {\vspace{-2pt}\end{shaded}\vspace{2pt}}

% Green: exploration / worked example / geometric reading
\newenvironment{greenbox}[1][Exploration for Self-Discovery]{%
  \vspace{4pt}\boxtitle{BoxGreenBd}{#1}%
  \colorlet{shadecolor}{BoxGreen}\begin{shaded}\vspace{-2pt}}%
  {\vspace{-2pt}\end{shaded}\vspace{2pt}}

\newenvironment{workedex}[1][Worked Example]{\begin{greenbox}[#1]}{\end{greenbox}}
\newenvironment{georeading}[1][Geometric \& Physical Reading]{\begin{greenbox}[#1]}{\end{greenbox}}

% Purple: TOK
\newenvironment{tokbox}[1][Theory of Knowledge]{%
  \vspace{4pt}\boxtitle{BoxPurpleBd}{#1}%
  \colorlet{shadecolor}{BoxPurple}\begin{shaded}\vspace{-2pt}}%
  {\vspace{-2pt}\end{shaded}\vspace{2pt}}

% Gray: info / note / international mindedness
\newenvironment{infobox}[1][Note]{%
  \vspace{4pt}\boxtitle{CMDark!70}{#1}%
  \colorlet{shadecolor}{CMGray}\begin{shaded}\vspace{-2pt}}%
  {\vspace{-2pt}\end{shaded}\vspace{2pt}}

% Red: broken math critique
\newenvironment{brokenmath}[1][The ``Broken Math'' Critique]{%
  \vspace{4pt}\boxtitle{BoxRedBd}{#1}%
  \colorlet{shadecolor}{BoxRed!50}\begin{shaded}\vspace{-2pt}}%
  {\vspace{-2pt}\end{shaded}\vspace{2pt}}

% ── Question environment ──────────────────────────────────────────────────────
\newcounter{qnum}
\newenvironment{question}[3]{%
  \refstepcounter{qnum}%
  \ifnum#1=1 \def\qshade{BoxBlue!20}\def\qtc{CMBlue}\def\qtier{$\bigstar$}%
  \else\ifnum#1=2 \def\qshade{BoxAmber!35}\def\qtc{CMAmber}\def\qtier{$\bigstar\bigstar$}%
  \else \def\qshade{BoxRed!20}\def\qtc{CMRed}\def\qtier{$\bigstar\bigstar\bigstar$}%
  \fi\fi
  \par\vspace{3pt}%
  \colorlet{shadecolor}{\qshade}%
  \begin{shaded}%
  \noindent\textbf{#3}\hfill\textcolor{\qtc}{\qtier}\par\smallskip}%
  {\end{shaded}\vspace{2pt}}

% ── Tearoff rule ──────────────────────────────────────────────────────────────
\newcommand{\tearoff}{\noindent\rule{\linewidth}{1.2pt}\par}

% ── Demand scale ──────────────────────────────────────────────────────────────
\newcommand{\demandscale}{%
\vspace{3pt}\noindent\small
\textbf{Demand scale (low $\to$ high):}\quad
\textcolor{CMBlue}{Write down} $<$
\textcolor{CMBlue}{State} $<$
\textcolor{CMBlue}{Describe} $<$
\textcolor{CMBlue}{Explain} $<$
\textcolor{CMBlue}{Determine} $<$
\textcolor{CMBlue}{Deduce} $<$
\textcolor{CMBlue}{Show that} $<$
\textcolor{CMRed}{\textbf{Prove}}%
\par\vspace{3pt}}

% ── Progress tracker ──────────────────────────────────────────────────────────
\newcommand{\progresstracker}{%
\noindent\small\textbf{Progress:}\enspace
P0~$\square$\enspace P1~$\square$\enspace P2~$\square$\enspace
P3~$\square$\enspace P4~$\square$\enspace P5~$\square$\enspace
P6~$\square$\enspace P7~$\square$\enspace P8~$\square$%
\par}

% ── Compulsory core ───────────────────────────────────────────────────────────
\newcommand{\compulsorycore}[1]{%
\vspace{2pt}%
\noindent\colorbox{BoxBlueBd}{\parbox{\linewidth}{\vspace{2pt}\hspace{5pt}%
  {\small\bfseries\color{white}Compulsory core
  ($\bigstar$ and $\bigstar\bigstar$):}\enspace
  {\small\color{white}Questions #1. \quad
  All $\bigstar\bigstar\bigstar$ questions are optional extensions.}%
  \vspace{2pt}}}\par\vspace{3pt}}

% ── Document metadata ─────────────────────────────────────────────────────────
\newcommand{\NATitle}{[Packet Title]}
\newcommand{\NASubtitle}{[Subtitle / Unifying Idea]}
\newcommand{\NACourse}{IBDP Mathematics AA HL}
\newcommand{\NATopics}{Topic [X]: [Name] \& Topic [Y]: [Name]}
\newcommand{\NAPrereqs}{[Prior activity names]}
\newcommand{\NAMaterials}{TI-84 Plus CE (or equivalent GDC).
Parts~0--5: Paper~1 style (non-calculator). GDC permitted in Part~6 only.}

% ── Title header ──────────────────────────────────────────────────────────────
\newcommand{\naheader}{%
  \noindent\colorbox{CMBlue}{\parbox{\linewidth}{%
    \vspace{3pt}\hspace{5mm}{\large\bfseries\color{white}CleverMathematics}%
    \hfill{\small\color{white!90!CMBlue}\NACourse}\hspace{5mm}\vspace{3pt}}}%
  \vspace{4mm}\par
  \noindent\begin{tabular}{@{}p{11cm}r@{}}
    \textbf{Student Name:}\quad\underline{\hspace{6cm}} &
    \textbf{Date:}\quad\underline{\hspace{3cm}}
  \end{tabular}\\[1mm]
  \noindent\textbf{Course:} \NACourse\\
  \noindent\textbf{Syllabus Topic(s):} \NATopics
  \vspace{3mm}\hrule\vspace{4mm}
  \begin{center}
    {\LARGE\bfseries\NATitle}\\[3pt]
    {\large\itshape\NASubtitle}
  \end{center}
  \vspace{2mm}
  \noindent\textbf{Materials:} \NAMaterials\\[1mm]
  \noindent\textbf{Prerequisites:} \NAPrereqs
  \vspace{2mm}
  \progresstracker
  \vspace{2mm}\hrule\vspace{4mm}}

% ── Part heading ──────────────────────────────────────────────────────────────
\newcounter{partnum}
\newcommand{\napart}[2]{%
  \refstepcounter{partnum}\label{part:#2}%
  \vspace{5pt}%
  \noindent\colorbox{CMBlue}{\parbox{\linewidth}{%
    \vspace{2pt}\hspace{5pt}{\large\bfseries\color{white}Part~\thepartnum: #1}\vspace{2pt}}}%
  \vspace{4pt}\par}

% ── Axis grid ──────────────────────────────────────────────────────────────────
% \axesgrid{xmin}{xmax}{ymin}{ymax}{width_cm}{height_cm}
\newcommand{\axesgrid}[6]{%
\begin{tikzpicture}[x=1cm,y=1cm]
  \draw[gray!30] (0,0) rectangle (#5,#6);
  \draw[gray!18,thin] (#5*0.2,0)--(#5*0.2,#6);
  \draw[gray!18,thin] (#5*0.4,0)--(#5*0.4,#6);
  \draw[gray!18,thin] (#5*0.6,0)--(#5*0.6,#6);
  \draw[gray!18,thin] (#5*0.8,0)--(#5*0.8,#6);
  \draw[gray!18,thin] (0,#6*0.2)--(#5,#6*0.2);
  \draw[gray!18,thin] (0,#6*0.4)--(#5,#6*0.4);
  \draw[gray!18,thin] (0,#6*0.6)--(#5,#6*0.6);
  \draw[gray!18,thin] (0,#6*0.8)--(#5,#6*0.8);
  \draw[-{Stealth[scale=0.7]},thick](0,#6*0.5)--(#5,#6*0.5)
      node[right,font=\footnotesize]{$x$};
  \draw[-{Stealth[scale=0.7]},thick](#5*0.5,0)--(#5*0.5,#6)
      node[above,font=\footnotesize]{$y$};
  \node[below,font=\tiny,gray] at (0,#6*0.5) {$#1$};
  \node[below,font=\tiny,gray] at (#5,#6*0.5) {$#2$};
  \node[left,font=\tiny,gray]  at (#5*0.5,0) {$#3$};
  \node[left,font=\tiny,gray]  at (#5*0.5,#6) {$#4$};
\end{tikzpicture}}

% ── Translation Table ─────────────────────────────────────────────────────────
% Usage:  \translationtable{Title}{Left Col}{Right Col}
%           row & row \\ row & row \\
%         \closett
\newcommand{\translationtable}[3]{%
  \vspace{4pt}%
  \noindent\colorbox{CMDark!25}{\parbox{\linewidth}{\vspace{1pt}\hspace{5pt}%
    {\small\bfseries Translation Table: #1}\vspace{1pt}}}\par\nobreak
  \colorlet{shadecolor}{CMGray}%
  \begin{shaded}\vspace{-4pt}%
  \begin{tabular}{@{}p{0.46\linewidth}@{\hspace{3pt}\vrule\hspace{3pt}}p{0.46\linewidth}@{}}%
  \rowcolor{CMDark!12}\small\textbf{#2} & \small\textbf{#3}\\
  \hline}
\newcommand{\closett}{\end{tabular}\vspace{-4pt}\end{shaded}\vspace{2pt}}

% ── Induction template ────────────────────────────────────────────────────────
\newcommand{\inductiontemplate}{%
  \vspace{3pt}%
  \colorlet{shadecolor}{BoxRed!20}%
  \begin{shaded}%
  \textbf{Proof by Mathematical Induction}\\[4pt]
  \textbf{Claim:}\hrulefill\\[10pt]
  \textbf{Base case} ($n=1$):\\[16pt]
  \rule{\linewidth}{0.2pt}\\[3pt]
  \textbf{Inductive hypothesis:} Assume the statement holds for $n=k$:\\[16pt]
  \rule{\linewidth}{0.2pt}\\[3pt]
  \textbf{Inductive step} (show true for $n=k+1$):\\[28pt]
  \rule{\linewidth}{0.2pt}\\[3pt]
  \textbf{Conclusion:} By the principle of mathematical induction, the statement
  holds for all $n\in\mathbb{Z}^+$.\hfill$\square$
  \end{shaded}\vspace{2pt}}

% ── Page style ────────────────────────────────────────────────────────────────
\pagestyle{fancy}
\fancyhf{}
\renewcommand{\headrulewidth}{0.4pt}
\fancyhead[L]{\small\itshape Nuanced Analysis: \NATitle}
\fancyhead[R]{\small\itshape CleverMathematics}
\fancyfoot[C]{\small\thepage}

% ============================================================
\begin{document}

\naheader

% ── Tear-off Command Terms Strip ──────────────────────────────────────────────
\tearoff
\boxtitle{BoxBlueBd}{Command Terms --- Tear-off Reference Strip}
\colorlet{shadecolor}{BoxBlue}
\begin{shaded}\small
\begin{tabular}{@{}>{\bfseries\color{CMBlue}}p{2.2cm}p{4.6cm}
                   >{\bfseries\color{CMBlue}}p{2.2cm}p{4.6cm}@{}}
Write down & No working required. A short answer.
  & Describe & Give a detailed account or picture.\\[1pt]
State      & Give a specific answer; no justification needed.
  & Explain  & Detailed account \emph{including reasons or causes.}\\[1pt]
Determine  & Obtain the only possible answer. Show your method.
  & Show that & Every step must appear. No shortcuts.\\[1pt]
Deduce     & Reach a conclusion from logical reasoning.
  & Prove    & Complete, rigorous chain of reasoning.\\[1pt]
Sketch     & Diagram with key features; relative scale required.
  & Hence    & \textbf{Must} use the immediately preceding result.\\[1pt]
Evaluate   & Find a numerical value or expression.
  & Hence or otherwise & Use previous result or any valid method.\\
\end{tabular}
\vspace{2pt}
\demandscale
\end{shaded}

\begin{ctbox}[Command-Term Spotlight: [TERM A] vs.\ [TERM B]]
\textbf{[TERM A]:} [Definition and what it demands from the student.]\\
\textbf{[TERM B]:} [Definition and what it demands from the student.]\\[2pt]
\textit{Key distinction:} [One sentence explaining when to use each and what the examiner
expects differently from each command in this packet.]
\end{ctbox}

\tearoff

\noindent\textbf{Vocabulary:}
\textbf{[Term 1]}, \textbf{[Term 2]}, \textbf{[Term 3]}, \textbf{[Term 4]},
\textbf{[Term 5]}, \textbf{[Term 6]}, \textbf{[Term 7]}.

\medskip
\noindent\textbf{ATLs:} \textit{Thinking \& Transfer Skills --- [One sentence naming the
Approaches to Learning skill built across this packet.]}

\begin{tokbox}[TOK Provocations --- Return to these in the Reflection (Q19)]
\begin{enumerate}
\item \textbf{[TOK Question 1.]} [Answerable using a specific result proven in this packet.]
\item \textbf{[TOK Question 2.]} [About discovery vs.\ invention, aesthetic beauty, or cultural attribution.]
\end{enumerate}
\end{tokbox}

\begin{infobox}[International Mindedness]
[Historical attribution. Include at least one non-European mathematician where genuinely
connected. Example: ``The series here were anticipated by \textbf{M\={a}dhava of
Sangam\={a}gr\={a}ma} (Kerala, India, c.\ 1400) --- roughly 250 years before Newton.'']
\end{infobox}

\compulsorycore{[Q1, Q2(a), Q2(b), Q3, Q4, Q5, Q6, Q8, Q9(a), Q9(b), Q10, Q12, Q13, Q15, Q17]}

\hrule

% ============================================================
%  PART 0: ACTIVATING PRIOR KNOWLEDGE
% ============================================================
\napart{Activating Prior Knowledge}{0}

\begin{microbox}
\begin{itemize}
  \item [Key fact or formula from the previous activity.]
  \item [Key formula: e.g., $z = r(\cos\theta + i\sin\theta)$, where $r=|z|$ and $\theta=\arg(z)$.]
  \item [Key concept 3.]
  \item [Key concept 4.]
\end{itemize}
\end{microbox}

\begin{question}{1}{1}{Q1 --- Write down}
\textbf{Write down} the [object] in [required form] for each of the following.\\[2pt]
(a)~$[z_1 = \ldots]$\quad(b)~$[z_2 = \ldots]$\quad(c)~$[z_3 = \ldots]$

\vspace{18pt}
\noindent
(a)\underline{\hspace{4.2cm}}\quad
(b)\underline{\hspace{4.2cm}}\quad
(c)\underline{\hspace{4.2cm}}
\end{question}

\begin{question}{1}{2}{Q2 --- Describe}
\textbf{Describe}, in one sentence each, what [the two quantities] tell you geometrically.
\textit{(You may answer with an annotated diagram.)}\\[2pt]
\textit{Sentence starter:} ``The [quantity] is [geometric meaning]\ldots''
\begin{enumerate}
  \item \underline{\hspace{\linewidth}}\\[10pt]\underline{\hspace{\linewidth}}
  \item \underline{\hspace{\linewidth}}\\[10pt]\underline{\hspace{\linewidth}}
\end{enumerate}
\end{question}

\begin{georeading}
[1--2 sentences translating Part~0's algebra into spatial language. Example: ``The modulus
$r$ gives the distance from the origin; $\theta$ gives the anticlockwise angle from the
positive real axis.'']
\end{georeading}

% ============================================================
%  PART 1: [CONJECTURE PHASE]
% ============================================================
\napart{[Part 1 Title --- Conjecture Phase]}{1}

\begin{microbox}
\begin{itemize}
  \item [Prerequisite A.]
  \item [Key formula from formula booklet: page X.]
  \item [Prerequisite C.]
\end{itemize}
\end{microbox}

\begin{question}{1}{3}{Q3 --- Numerical Warm-up}
[Context sentence.] Let $[z_1 = \text{specific value}]$ and $[z_2 = \text{specific value}]$.

\textbf{Compute} [the product/result] and \textbf{write down} its [property A] and [property B].

\vspace{12pt}
\noindent[Property A]${}=\underline{\hspace{2.8cm}}$\hfill
[Property B]${}=\underline{\hspace{2.8cm}}$
\end{question}

\begin{question}{2}{4}{Q4 --- Conjecture}
Based on Q3, \textbf{write} a conjecture relating [the result] to [its components].

\textit{Sentence starter:} ``I conjecture that\ldots''

\vspace{34pt}
\end{question}

\begin{question}{2}{5}{Q5 --- Show that}
\textbf{Show that} your conjecture from Q4 holds in general by [method], using
the \textbf{[formula name]} from the formula booklet.

\begin{amberbox}[First line given --- working required from here]
$[z_1\cdot z_2] = [r_1][r_2](\cos\alpha+i\sin\alpha)(\cos\beta+i\sin\beta)$
\end{amberbox}

\vspace{52pt}
\end{question}

\begin{georeading}
[Geometric interpretation. Example: ``Multiplying two complex numbers \emph{rotates} by
the sum of arguments and \emph{scales} by the product of moduli.'']
\end{georeading}

\translationtable{[Domain Transfer Name: e.g., Roots and Multiplicity]}
{What you observe on screen\ldots}{What you write on the exam paper\ldots}
[Observation 1] & [Formal IB phrasing 1]\\[3pt]
[Observation 2] & [Formal IB phrasing 2]\\[3pt]
[Observation 3] & [Formal IB phrasing 3]\\
\closett

% ============================================================
%  PART 2: [PROOF PHASE]
% ============================================================
\napart{[Part 2 Title --- Proof Phase]}{2}

\begin{microbox}
\begin{itemize}
  \item [Key result from Part~1.]
  \item [Formula booklet reference: page X.]
\end{itemize}
\end{microbox}

\begin{question}{1}{6}{Q6 --- Numerical Investigation}
Restrict to [the unit circle / a specific case]. Using your Part~1 result repeatedly,
\textbf{write} a conjecture for $[\text{expression}]^n$ where $n\in\mathbb{Z}^+$.

\vspace{34pt}
\end{question}

\begin{question}{3}{7}{Q7 --- Prove by Mathematical Induction}
\textbf{Prove} that for all $n\in\mathbb{Z}^+$,
\[
[\text{claim, e.g., }(\cos\theta+i\sin\theta)^n = \cos n\theta + i\sin n\theta]
\]
\inductiontemplate
\end{question}

\begin{question}{2}{8}{Q8 --- Deduce}
\textbf{Deduce} (using Q7) [the extended general version of the result].

\textit{Hint:} [One-sentence minimal hint.]

\vspace{34pt}
\end{question}

\begin{brokenmath}
\textit{The following working was submitted by a student. Your job is not to judge
the student --- errors like this reveal important mathematical distinctions. Find
the slip and explain its consequence.}

\bigskip
\noindent\textbf{The problem:} [State the problem clearly.]

\medskip
\noindent\textbf{Student's working:}
\begin{align*}
  [\text{Step 1}] &= [\text{correct intermediate result}]\\
  [\text{Step 2}] &= [\text{result with FATAL ERROR}]\quad\text{[student's comment]}\\
  [\text{Step 3}] &= [\text{wrong final answer}]
\end{align*}
\end{brokenmath}

\begin{question}{2}{9}{Q9(a) --- Explain}
\textbf{Explain} the fatal conceptual error. \textit{(You may answer in bullet points.)}

\textit{Sentence starter:} ``This error occurs because\ldots, which means that\ldots''

\vspace{28pt}
\end{question}

\begin{question}{2}{9}{Q9(b) --- Determine}
\textbf{Determine} the correct value and state it in [required form].

\vspace{28pt}
\end{question}

\begin{georeading}
[Geometric/physical interpretation of Part~2's main result. 1--2 sentences.]
\end{georeading}

% ============================================================
%  PART 3: [APPLICATION PHASE]
% ============================================================
\napart{[Part 3 Title --- Application Phase]}{3}

\begin{microbox}
\begin{itemize}
  \item [Key result from Part~2.]
  \item [Formula booklet reference: page X.]
\end{itemize}
\end{microbox}

\begin{question}{2}{10}{Q10 --- Worked Example then Practice}
\begin{workedex}
[Simpler analogous problem fully worked, one level below the target difficulty.]\par
\textbf{Find} [simpler case, e.g., $n=2$]:
\begin{align*}
  [\text{Step 1}] &= [\text{intermediate result}]\\
  [\text{Step 2}] &= [\text{final result}]
\end{align*}
\end{workedex}

\textbf{Hence} (using the result above), \textbf{find} [the harder full case].

\vspace{40pt}
\end{question}

\begin{question}{2}{11}{Q11 --- Show that}
\textbf{Show that} [second result at full difficulty; no model provided].

\vspace{40pt}
\end{question}

\begin{question}{1}{12}{Q12 --- Sketch}
\textbf{Sketch} $y=[f(x)]$ on the axes provided. Label all intercepts, extrema, inflection points.
\textit{(You may attach an annotated GeoGebra screenshot.)}

\begin{center}
\axesgrid{-5}{5}{-4}{4}{9}{7}
\end{center}
\end{question}

% ============================================================
%  PART 4: [CALCULUS / ANALYSIS PHASE]
% ============================================================
\napart{[Part 4 Title --- Calculus/Analysis Phase]}{4}

\begin{microbox}
\begin{itemize}
  \item [Key algebraic result from Part~3.]
  \item $\displaystyle f'(x)=\lim_{h\to 0}\dfrac{f(x+h)-f(x)}{h}$
  \item [Formula booklet reference: page X.]
\end{itemize}
\end{microbox}

\begin{question}{2}{13}{Q13 --- Numerical Estimate}
Consider $f(x)=[g(x)]$. Find the [slope/derivative] at $x=[a]$.
\begin{enumerate}
  \item \textbf{Evaluate} $f([a])$ and $f([a]+0.1)$. Calculate the secant slope.
  \vspace{16pt}
  \item \textbf{Determine} a simplified expression for the secant slope from $x=[a]$ to $x=[a]+\Delta x$.
  \vspace{30pt}
  \item \textbf{Evaluate} $\displaystyle\lim_{\Delta x\to 0}$ [your expression from (b)].
  \vspace{16pt}
\end{enumerate}
\end{question}

\begin{question}{3}{14}{Q14 --- First Principles Proof}
\textbf{Show that} the derivative of $f(x)=[g(x)]$ is $f'(x)=[g'(x)]$:
\[f'(x)=\lim_{h\to 0}\frac{f(x+h)-f(x)}{h}\]

\begin{amberbox}[Scaffold --- fold under if not needed]
Expand using [Binomial Theorem]. The $f(x)$ terms cancel. Factor $h$, then evaluate.
\end{amberbox}

\vspace{52pt}
\end{question}

\begin{georeading}
[$f'(x)$ gives the tangent gradient at any point $x$. At a local extremum: $f'(x)=0$.]
\end{georeading}

\translationtable{Calculus Features in Context}
{Algebraic feature}{Applied context}
$x$-intercepts ($y=0$)  & ``Break-even'', ``Hits the ground'', ``Sea level''\\[3pt]
$f'(x)=0$ (turning point) & ``Max profit'', ``Particle at rest'', ``Peak height''\\[3pt]
$f''(x)=0$ (inflection)   & ``Diminishing returns'', ``Max acceleration''\\[3pt]
[Feature 4] & [Context 4]\\
\closett

% ============================================================
%  PART 5: [SYNTHESIS PHASE]
% ============================================================
\napart{[Part 5 Title --- Synthesis Phase]}{5}

\begin{microbox}
\begin{itemize}
  \item [Key result from Part~4.]
  \item [How algebraic and geometric routes converge.]
\end{itemize}
\end{microbox}

\begin{question}{2}{15}{Q15 --- Two Representations}
\textbf{Explain} why [Methods A and B] give the same result.
\textit{(You may answer in bullet points.)}

\textit{Sentence starter:} ``Both methods agree because\ldots''

\vspace{34pt}
\end{question}

\begin{question}{3}{16}{Q16 --- Third Representation}
[Optional $\bigstar\bigstar\bigstar$. Push the same object into a third form.]

\textbf{Determine} [third representation or application].

\vspace{40pt}
\end{question}

% ============================================================
%  PART 6: GDC MASTERY
% ============================================================
\napart{GDC Mastery}{6}\label{part:gdc}

\begin{microbox}
\begin{itemize}
  \item [Calculus result from Part~4 for window optimisation.]
  \item GDC permitted throughout this Part.
  \item Store your function as \texttt{Y1}.
\end{itemize}
\end{microbox}

\begin{infobox}[GDC Efficiency: The Window Wars]
Use calculus to determine the optimal viewing window \emph{before} pressing \textsc{Graph}.
Store complex functions as \texttt{Y1} so you never retype them.
\end{infobox}

\begin{question}{2}{17}{Q17 --- Optimising the GDC Window}
Consider $f(x)=[expression]$. A standard window shows nothing.
\begin{enumerate}
  \item \textbf{Determine} $f'(x)$ and the exact $x$-coordinates of all local extrema.
  \vspace{16pt}
  \item \textbf{Evaluate} $y$-values at each extremum. \textbf{State} the required $y$-range.
  \vspace{16pt}
  \item \textbf{State} an appropriate window:
  \vspace{4pt}\\
  \begin{tabular}{@{}ll@{\qquad}ll@{}}
    Xmin${}=\underline{\hspace{2cm}}$ & Xmax${}=\underline{\hspace{2cm}}$ &
    Ymin${}=\underline{\hspace{2cm}}$ & Ymax${}=\underline{\hspace{2cm}}$
  \end{tabular}
\end{enumerate}
\end{question}

\begin{question}{1}{18}{Q18 --- GDC Sketch}
Using your window from Q17(c), \textbf{sketch} $f(x)$ below.
\textit{(You may attach a GeoGebra screenshot.)}

\begin{center}
\axesgrid{-2}{15}{-6}{4}{10}{7}
\end{center}
\end{question}

% ============================================================
%  PART 7: REFLECTION
% ============================================================
\napart{Reflection (Metacognition)}{7}

\noindent\textit{You may respond orally --- ask your teacher to record a voice memo.
Bullet points are acceptable throughout.}

\begin{question}{1}{19}{Q19 --- Concept Map}
\textbf{List} the major concepts and connections. Aim for six; one from each syllabus topic.

\vspace{5pt}
\begin{tabular}{@{}>{\centering\arraybackslash}p{3cm}
                    >{\centering\arraybackslash}p{4cm}
                    >{\centering\arraybackslash}p{7cm}@{}}
\rowcolor{CMGray}
\small\textbf{Concept} & \small\textbf{Where it appeared} &
\small\textbf{How it connected to another concept}\\
\hline
& &\\[9pt]\hline
& &\\[9pt]\hline
& &\\[9pt]\hline
& &\\[9pt]\hline
& &\\[9pt]\hline
& &\\[9pt]\hline
\end{tabular}
\end{question}

\begin{question}{2}{20}{Q20 --- Multiple Representations}
\textbf{Explain} what is gained by two independent proofs of one result.
Does a second proof make the result \emph{more true}?
\textit{(You may answer in bullet points.)}

\textit{Sentence starter:} ``A second proof is valuable because\ldots''

\vspace{40pt}
\end{question}

\begin{question}{3}{21}{Q21 --- TOK Position Statement}
Return to \textbf{TOK Provocation~1} and \textbf{take and defend a position}.

\vspace{4pt}
\colorlet{shadecolor}{BoxPurple!60}
\begin{shaded}\itshape\small
``I argue that \underline{\hspace{5.5cm}}. My evidence from this packet is
\underline{\hspace{4.5cm}}. A counterargument would be \underline{\hspace{4cm}},
but I respond that \underline{\hspace{5cm}}.''
\end{shaded}

\vspace{3pt}
\noindent\textit{Modelled mentor text (different result --- read, then write your own):}\\
\textit{``[One example paragraph using a result the student will NOT choose, showing the genre.]''}

\vspace{32pt}
\end{question}

% ============================================================
%  EXTENSION & IA-SEEDING
% ============================================================
\napart{Extension \& IA-Seeding --- \textit{Optional: choose one}}{8}

\begin{infobox}[Note for Students]
These branches are deliberately \emph{under-specified}, in the spirit of the IB Internal
Assessment. There is no single correct approach.
\end{infobox}

\begin{question}{3}{E1}{Branch E1 --- [Topic Area A]}
\textbf{[Opening sentence connecting to a different IB topic area.]}

[2--3 sentences. Keep deliberately open-ended.]

\textit{IA relevance:} [One sentence.]
\end{question}

\begin{question}{3}{E2}{Branch E2 --- [Topic Area B]}
\textbf{[Opening sentence connecting to a second IB topic area.]}

[2--3 sentences.]

\textit{IA relevance:} [One sentence.]
\end{question}

% ============================================================
%  TEACHER'S COMPANION  (Remove before distributing to students)
% ============================================================
\newpage
\noindent\colorbox{CMBlue}{\parbox{\linewidth}{%
  \vspace{4pt}\hspace{6pt}{\Large\bfseries\color{white}Teacher's Companion}\\[2pt]
  \hspace{6pt}{\small\itshape\color{white!80!CMBlue}Remove before distributing. Instructor use only.}%
  \vspace{4pt}}}
\vspace{5pt}

\subsection*{A. Integration Map}
\begin{tabular}{@{}>{\bfseries}p{4cm}p{10.5cm}@{}}
\toprule
IB Element & Location in packet\\
\midrule
Topic [X] & [Q numbers]\\
Topic [Y] & [Q numbers]\\
TOK & [Header + Q21]\\
ATL & [Q15, Q20, Q21]\\
International Mindedness & [IM box + Q[N]]\\
GDC Technology & [Q17--Q18]\\
IA Seeding & [Branches E1, E2]\\
Paper 1 & [Parts~0--5]\\
Paper 3 & [Parts~1--4]\\
Command terms & [Write down, Describe, Explain, Determine, Deduce, Show that, Prove, Hence, Sketch]\\
\bottomrule
\end{tabular}

\subsection*{B. Model's Moves}
\begin{itemize}
  \item \textbf{Conjecture-before-rule:} Q4$\to$Q5 and Q6$\to$Q7.
  \item \textbf{Numerical warm-up:} Q3 (specific) $\to$ Q4--5 (general).
  \item \textbf{Worked example $\to$ practice:} Q10 $\to$ Q11.
  \item \textbf{Planted error:} Q9 ([Misconception name]).
  \item \textbf{Translation tables:} Part~1 and Part~4.
  \item \textbf{Rule of four:} algebraic (Q5), geometric (Q12), numerical (Q3), verbal (Q15).
  \item \textbf{Proof scaffolds:} Q7 (induction template), Q14 (first line given).
  \item \textbf{GDC window from calculus:} Q17--Q18.
  \item \textbf{Metacognition:} Q19, Q20, Q21.
  \item \textbf{IA seeding:} E1, E2 (under-specified).
\end{itemize}

\subsection*{C. Answer Sketches \& Planted-Error Key}
\begin{itemize}
  \item \textbf{Q1:} [(a) (b) (c)]
  \item \textbf{Q2:} (a)~[Answer] (b)~[Answer]
  \item \textbf{Q3:} [Property A]$=[\ldots]$; [Property B]$=[\ldots]$
  \item \textbf{Q4:} [Expected conjecture]
  \item \textbf{Q5:} [Proof outline]
  \item \textbf{Q6:} [Conjecture]
  \item \textbf{Q7:} [Base case: $n=1$ straightforward. Inductive step: applies Q5 with [substitution]. Conclusion.]
  \item \textbf{Q8:} [Deduction: $r^n$ carries through.]
  \item \textbf{Q9:} Correct: $[\ldots]$. Misconception: ``[Name]''. HL concept: [distinguished idea].
  \item \textbf{Q10--Q11:} [Answers]
  \item \textbf{Q12:} [Intercepts, extrema, end behaviour]
  \item \textbf{Q13:} (a)~$[\ldots]$ (b)~$[\ldots]$ (c)~$[\ldots]$
  \item \textbf{Q14:} [First-principles outline]
  \item \textbf{Q17:} $f'(x)=[\ldots]$; extrema $x=[\ldots]$; window $[\ldots]$
\end{itemize}

\subsection*{D. Tiered Deadline Guidance}
\begin{tabular}{@{}p{3.8cm}p{10.8cm}@{}}
\toprule
\textbf{Session} & \textbf{Parts}\\
\midrule
50-min lesson & Parts~0--2 (Q1--Q9).\\
Double period & Parts~0--4 (Q1--Q14).\\
Take-home & Full packet + one extension branch.\\
Reduced workload & Compulsory core only.\\
\bottomrule
\end{tabular}

\subsection*{E. Compulsory Core}
Q1, Q2, Q3, Q4, Q5, Q6, Q8, Q9(a), Q9(b), Q10, Q12, Q13, Q17.

\subsection*{F. Differentiation Note}
\begin{tabular}{@{}p{3.8cm}p{10.8cm}@{}}
\toprule
\textbf{Profile} & \textbf{Guidance}\\
\midrule
ELL & Tear-off strip essential. Sentence starters reduce language demand. Accept diagrams.\\[4pt]
ADHD/dyslexia & Progress tracker front page. Frame Q9 positively.\\[4pt]
Gaps & Verify Part~0 first. Run [prior activity] if needed.\\[4pt]
Gifted & Hide scaffolds. Direct to $\bigstar\bigstar\bigstar$ and IA branches.\\
\bottomrule
\end{tabular}

\subsection*{G. Design Note}
\textbf{Genuinely integrated:} [Topic X] and [Topic Y] --- [explain the mathematical thread].

\textbf{Extension-level only:} [Topic Z] --- [honest reason why it isn't in the core].

\end{document}
